% Source: https://artofproblemsolving.com/community/c3h1931055_binary_numbers_and_polygons
% Stable identity: aops-topic-1931055-q01
The integers $0$ through $63$ are exactly the numbers with at most $6$ binary digits. Pad each with leading zeros to six digits (this adds no $1$'s); then the list consists of all $2^6=64$ strings of six binary digits. In each of the six positions, exactly half of the strings, $32$, have a $1$. So the number of $1$'s is $6\cdot32=\boxed{192}$. (The same count for $0$ through $7$ gives $3\cdot4=12$, as stated.)
